\section{The First Permission}

Measurement cannot begin by saying what a thing is. To say what a
thing is already presumes a stable place from which the saying can be
checked. It presumes a sameness between the thing indicated now and the
thing indicated again, a sameness between the mark and what the mark is
allowed to carry, and a sameness between the reader before and after the
act of reading. Those presumptions are too expensive for a foundation.
At the beginning they are not yet available.

The grammar therefore begins with a weaker permission. It permits a
difference to be read. It does not yet identify the two sides of the
difference as objects, does not yet count them, and does not yet decide
whether either side is true. It only permits the first separation: this
is not that. The separation is not a theory of the world. It is the
minimal condition under which a theory could later be answerable to the
world.

This chapter follows the consequences of that permission. Once
distinguishability is allowed, counting becomes possible, but counting
already splits into two readings. A reading can preserve a whole, or it
can break the whole into parts. Once parts can be carried, a number no
longer needs to arrive as an infinite object already complete. It can
arrive as a finite condition with a residue still to be reduced. Once
residue can be reduced, a clock enters: not as a metaphysical river, but
as a permitted alternation by which a trial can be repeated. Once trials
are repeated, observation is bounded by the time that can be spent.

The bound is not a defect added from outside. It is what keeps the
ledger honest. A measurement that could spend no time would have no
trial; a measurement that could spend all time would have no finite
record. The grammar must therefore carry a horizon. It may refine, but
it may not pretend that refinement has already crossed every possible
boundary.

The danger appears when the grammar turns back on itself. The tower
built from distinction, counting, finite approach, clock, trial, and
bound can be pointed at truth itself. At that bare foundation the grammar
finds no interior difference between affirmation and denial. This is the
collapse named in the chapter title. It is not a theatrical confusion of
ordinary yes and no. It is a disciplined failure: the grammar refuses to
use a distinction that it has not earned.

The chapter ends with the one act that can follow such a collapse. A
single sanctioned identity is manufactured. It is not installed as a new
primitive. It is a needle: local, bounded, and answerable to the
distinction from which the whole construction began. With that needle
the grammar restores a place to stand, and the rest of the book can ask
what finite measurement is forced to carry once truth has survived its
own collapse.
